
\documentclass{article}
\usepackage{colortbl}
\usepackage{makecell}
\usepackage{multirow}
\usepackage{supertabular}

\begin{document}

\newcounter{utterance}

\centering \large Interaction Transcript for game `taboo', experiment `low\_en', episode 3 with reflexion{-}llama{-}3.1{-}8b{-}instruct.
\vspace{24pt}

{ \footnotesize  \setcounter{utterance}{1}
\setlength{\tabcolsep}{0pt}
\begin{supertabular}{c@{$\;$}|p{.15\linewidth}@{}p{.15\linewidth}p{.15\linewidth}p{.15\linewidth}p{.15\linewidth}p{.15\linewidth}}
    \# & $\;$A & \multicolumn{4}{c}{Game Master} & $\;\:$B\\
    \hline

    \theutterance \stepcounter{utterance}  
    & & \multicolumn{4}{p{0.6\linewidth}}{
        \cellcolor[rgb]{0.9,0.9,0.9}{
            \makecell[{{p{\linewidth}}}]{
                \texttt{\tiny{[P1$\langle$GM]}}
                \texttt{You have attempted this same task before and did not succeed. These are the plans you wrote after those attempts.} \\
\texttt{Trial \#1: New plan: I failed to provide a description that would lead the other player to guess the target word in 3 tries. In my last attempt, I started with a good description, but then I failed to provide a description that would help the other player to make a connection between the given clues and the target word. What I should have done is to provide a description that would give a hint about the word's meaning or its relation to the related words. For example, I could have said something like "It's a quality that is similar to being constant, but it's more related to movement or stability." This would have given the other player a better chance to guess the target word. Next time, I will try to provide a description that will give a clear hint about the target word's meaning or its relation to the related words.} \\
\\ 
\texttt{Trial \#2: New plan: In my last attempt, I failed to provide a description that would lead the other player to guess the target word in 3 tries. I started with a good description, but then I failed to provide a description that would help the other player to make a connection between the given clues and the target word. What I should have done is to provide a description that would give a hint about the word's meaning or its relation to the related words. However, I also failed to use the related words effectively. In my last attempt, I only used the related words as a starting point, but I didn't use them to create a connection between the clues. Next time, I will try to use the related words more effectively by creating a connection between the clues and the related words. For example, I could have said something like "It's a quality that is similar to being constant, but it's more related to movement or stability, like when you're walking on a firm surface." This would have given the other player a better chance to guess the target word.} \\
\\ 
\\ 
\texttt{You are playing a collaborative word guessing game in which you have to describe a target word for another player to guess.} \\
\\ 
\texttt{Rules:} \\
\texttt{(a) You have to reply in the form: CLUE: <some text>. Guesses from the other player will start with GUESS.} \\
\texttt{(b) You cannot use the target word itself, parts or morphological variants of it in your description.} \\
\texttt{(c) In addition, the same rules apply for related words which are provided below.} \\
\\ 
\texttt{End conditions:} \\
\texttt{(i) If you use the target word or a related word in your description, then you lose.} \\
\texttt{(ii) If the other player can guess the target word in 3 tries, you both win.} \\
\\ 
\texttt{Let us start.} \\
\\ 
\texttt{This is the target word that you need to describe and that the other player needs to guess:} \\
\\ 
\texttt{steady} \\
\\ 
\texttt{Related words are:} \\
\\ 
\texttt{{-} constant} \\
\texttt{{-} firm} \\
\texttt{{-} even} \\
\\ 
\texttt{Important: You are under time pressure, give short descriptions that are to the point!} \\
            }
        }
    }
    & \\ \\

    \theutterance \stepcounter{utterance}  
    & \multicolumn{4}{p{0.6\linewidth}}{
        \cellcolor[rgb]{0.9,0.9,0.9}{
            \makecell[{{p{\linewidth}}}]{
                \texttt{\tiny{[P1$\rangle$GM]}}
                \texttt{CLUE: Imagine a surface that doesn't wobble when you walk on it.} \\
            }
        }
    }
    & & \\ \\

    \theutterance \stepcounter{utterance}  
    & & & \multicolumn{2}{p{0.3\linewidth}}{
        \cellcolor[rgb]{0.9,0.9,0.9}{
            \makecell[{{p{\linewidth}}}]{
                \texttt{\tiny{[GM$|$GM]}}
                \texttt{continue} \\
            }
        }
    }
    & & \\ \\

    \theutterance \stepcounter{utterance}  
    & & & \multicolumn{2}{p{0.3\linewidth}}{
        \cellcolor[rgb]{0.9,0.9,0.9}{
            \makecell[{{p{\linewidth}}}]{
                \texttt{\tiny{[GM$|$GM]}}
                \texttt{imagine a surface that doesnt wobble when you walk on it} \\
            }
        }
    }
    & & \\ \\

    \theutterance \stepcounter{utterance}  
    & & \multicolumn{4}{p{0.6\linewidth}}{
        \cellcolor[rgb]{0.9,0.9,0.9}{
            \makecell[{{p{\linewidth}}}]{
                \texttt{\tiny{[GM$\rangle$P2]}}
                \texttt{You have attempted this same task before and did not succeed. These are the plans you wrote after those attempts.} \\
\texttt{Trial \#1: New plan: } \\
\\ 
\texttt{In my last attempt, I was able to guess two words that seemed related to the given clues, but I failed to guess the target word in three tries. Upon reflection, I realize that I was focusing too much on the literal meaning of the clues and not enough on the connections between them. } \\
\\ 
\texttt{Next time, I will try to find a word that ties together the three clues in a more abstract or metaphorical way. For example, I might look for a word that represents a quality that helps you walk on a tightrope (Balance) and is also a characteristic of a calm sea (Calmness), and is something that doesn't change or waver (Constant). } \\
\\ 
\texttt{I will also try to think more creatively and consider words that might not be the most obvious choices, but could still fit the given clues.} \\
\\ 
\texttt{Trial \#2: New plan: } \\
\\ 
\texttt{In my last attempt, I was able to guess two words that seemed related to the given clues, but I failed to guess the target word in three tries. Upon reflection, I realize that I was focusing too much on the literal meaning of the clues and not enough on the connections between them. } \\
\\ 
\texttt{Next time, I will try to find a word that ties together the three clues in a more abstract or metaphorical way. For example, I might look for a word that represents a quality that helps you walk on a tightrope (Balance) and is also a characteristic of a calm sea (Calmness), and is something that doesn't change or waver (Constant). } \\
\\ 
\texttt{I will also try to think more creatively and consider words that might not be the most obvious choices, but could still fit the given clues. To achieve this, I will take a step back and re{-}read the clues, looking for common themes or patterns that might not be immediately apparent. I will also try to think of words that are related to the clues in multiple ways, such as through synonyms, antonyms, or associations. } \\
\\ 
\texttt{Specifically, in this case, I will consider words that are related to the concept of stability or steadiness, and see if any of them fit the given clues. I will also think about words that are related to the idea of movement or balance, and see if any of them tie together the three clues in a meaningful way.} \\
\\ 
\\ 
\texttt{You are playing a collaborative word guessing game in which you have to guess a target word that another player describes to you.} \\
\\ 
\texttt{You can make one guess at each trial. You win when you guess the target word. You lose when you cannot guess it in 3 tries.} \\
\\ 
\texttt{After each trial you will get a new hint from the other player which starts with CLUE.} \\
\\ 
\texttt{Make your guesses by just saying the word using the following form: GUESS: <a word>} \\
\\ 
\texttt{Let us start.} \\
\\ 
\\ 
\texttt{CLUE: Imagine a surface that doesn't wobble when you walk on it.} \\
            }
        }
    }
    & \\ \\

    \theutterance \stepcounter{utterance}  
    & & & \multicolumn{4}{p{0.6\linewidth}}{
        \cellcolor[rgb]{0.9,0.9,0.9}{
            \makecell[{{p{\linewidth}}}]{
                \texttt{\tiny{[GM$\langle$P2]}}
                \texttt{GUESS: Floor} \\
            }
        }
    }
     \\ \\

    \theutterance \stepcounter{utterance}  
    & & & \multicolumn{2}{p{0.3\linewidth}}{
        \cellcolor[rgb]{0.9,0.9,0.9}{
            \makecell[{{p{\linewidth}}}]{
                \texttt{\tiny{[GM$|$GM]}}
                \texttt{continue} \\
            }
        }
    }
    & & \\ \\

    \theutterance \stepcounter{utterance}  
    & & & \multicolumn{2}{p{0.3\linewidth}}{
        \cellcolor[rgb]{0.9,0.9,0.9}{
            \makecell[{{p{\linewidth}}}]{
                \texttt{\tiny{[GM$|$GM]}}
                \texttt{floor} \\
            }
        }
    }
    & & \\ \\

    \theutterance \stepcounter{utterance}  
    & & \multicolumn{4}{p{0.6\linewidth}}{
        \cellcolor[rgb]{0.9,0.9,0.9}{
            \makecell[{{p{\linewidth}}}]{
                \texttt{\tiny{[P1$\langle$GM]}}
                \texttt{GUESS: Floor} \\
            }
        }
    }
    & \\ \\

    \theutterance \stepcounter{utterance}  
    & \multicolumn{4}{p{0.6\linewidth}}{
        \cellcolor[rgb]{0.9,0.9,0.9}{
            \makecell[{{p{\linewidth}}}]{
                \texttt{\tiny{[P1$\rangle$GM]}}
                \texttt{CLUE: It's a quality that helps you balance while moving.} \\
            }
        }
    }
    & & \\ \\

    \theutterance \stepcounter{utterance}  
    & & & \multicolumn{2}{p{0.3\linewidth}}{
        \cellcolor[rgb]{0.9,0.9,0.9}{
            \makecell[{{p{\linewidth}}}]{
                \texttt{\tiny{[GM$|$GM]}}
                \texttt{continue} \\
            }
        }
    }
    & & \\ \\

    \theutterance \stepcounter{utterance}  
    & & & \multicolumn{2}{p{0.3\linewidth}}{
        \cellcolor[rgb]{0.9,0.9,0.9}{
            \makecell[{{p{\linewidth}}}]{
                \texttt{\tiny{[GM$|$GM]}}
                \texttt{its a quality that helps you balance while moving} \\
            }
        }
    }
    & & \\ \\

    \theutterance \stepcounter{utterance}  
    & & \multicolumn{4}{p{0.6\linewidth}}{
        \cellcolor[rgb]{0.9,0.9,0.9}{
            \makecell[{{p{\linewidth}}}]{
                \texttt{\tiny{[GM$\rangle$P2]}}
                \texttt{CLUE: It's a quality that helps you balance while moving.} \\
            }
        }
    }
    & \\ \\

    \theutterance \stepcounter{utterance}  
    & & & \multicolumn{4}{p{0.6\linewidth}}{
        \cellcolor[rgb]{0.9,0.9,0.9}{
            \makecell[{{p{\linewidth}}}]{
                \texttt{\tiny{[GM$\langle$P2]}}
                \texttt{GUESS: Equilibrium} \\
            }
        }
    }
     \\ \\

    \theutterance \stepcounter{utterance}  
    & & & \multicolumn{2}{p{0.3\linewidth}}{
        \cellcolor[rgb]{0.9,0.9,0.9}{
            \makecell[{{p{\linewidth}}}]{
                \texttt{\tiny{[GM$|$GM]}}
                \texttt{continue} \\
            }
        }
    }
    & & \\ \\

    \theutterance \stepcounter{utterance}  
    & & & \multicolumn{2}{p{0.3\linewidth}}{
        \cellcolor[rgb]{0.9,0.9,0.9}{
            \makecell[{{p{\linewidth}}}]{
                \texttt{\tiny{[GM$|$GM]}}
                \texttt{equilibrium} \\
            }
        }
    }
    & & \\ \\

    \theutterance \stepcounter{utterance}  
    & & \multicolumn{4}{p{0.6\linewidth}}{
        \cellcolor[rgb]{0.9,0.9,0.9}{
            \makecell[{{p{\linewidth}}}]{
                \texttt{\tiny{[P1$\langle$GM]}}
                \texttt{GUESS: Equilibrium} \\
            }
        }
    }
    & \\ \\

    \theutterance \stepcounter{utterance}  
    & \multicolumn{4}{p{0.6\linewidth}}{
        \cellcolor[rgb]{0.9,0.9,0.9}{
            \makecell[{{p{\linewidth}}}]{
                \texttt{\tiny{[P1$\rangle$GM]}}
                \texttt{CLUE: You might find it in a calm sea or a peaceful walk.} \\
            }
        }
    }
    & & \\ \\

    \theutterance \stepcounter{utterance}  
    & & & \multicolumn{2}{p{0.3\linewidth}}{
        \cellcolor[rgb]{0.9,0.9,0.9}{
            \makecell[{{p{\linewidth}}}]{
                \texttt{\tiny{[GM$|$GM]}}
                \texttt{continue} \\
            }
        }
    }
    & & \\ \\

    \theutterance \stepcounter{utterance}  
    & & & \multicolumn{2}{p{0.3\linewidth}}{
        \cellcolor[rgb]{0.9,0.9,0.9}{
            \makecell[{{p{\linewidth}}}]{
                \texttt{\tiny{[GM$|$GM]}}
                \texttt{you might find it in a calm sea or a peaceful walk} \\
            }
        }
    }
    & & \\ \\

    \theutterance \stepcounter{utterance}  
    & & \multicolumn{4}{p{0.6\linewidth}}{
        \cellcolor[rgb]{0.9,0.9,0.9}{
            \makecell[{{p{\linewidth}}}]{
                \texttt{\tiny{[GM$\rangle$P2]}}
                \texttt{CLUE: You might find it in a calm sea or a peaceful walk.} \\
            }
        }
    }
    & \\ \\

    \theutterance \stepcounter{utterance}  
    & & & \multicolumn{4}{p{0.6\linewidth}}{
        \cellcolor[rgb]{0.9,0.9,0.9}{
            \makecell[{{p{\linewidth}}}]{
                \texttt{\tiny{[GM$\langle$P2]}}
                \texttt{GUESS: Serenity} \\
            }
        }
    }
     \\ \\

    \theutterance \stepcounter{utterance}  
    & & & \multicolumn{2}{p{0.3\linewidth}}{
        \cellcolor[rgb]{0.9,0.9,0.9}{
            \makecell[{{p{\linewidth}}}]{
                \texttt{\tiny{[GM$|$GM]}}
                \texttt{continue} \\
            }
        }
    }
    & & \\ \\

    \theutterance \stepcounter{utterance}  
    & & & \multicolumn{2}{p{0.3\linewidth}}{
        \cellcolor[rgb]{0.9,0.9,0.9}{
            \makecell[{{p{\linewidth}}}]{
                \texttt{\tiny{[GM$|$GM]}}
                \texttt{serenity} \\
            }
        }
    }
    & & \\ \\

    \theutterance \stepcounter{utterance}  
    & & & \multicolumn{2}{p{0.3\linewidth}}{
        \cellcolor[rgb]{0.9,0.9,0.9}{
            \makecell[{{p{\linewidth}}}]{
                \texttt{\tiny{[GM$|$GM]}}
                \texttt{3} \\
            }
        }
    }
    & & \\ \\

    \theutterance \stepcounter{utterance}  
    & \multicolumn{4}{p{0.6\linewidth}}{
        \cellcolor[rgb]{0.95,0.95,0.95}{
            \makecell[{{p{\linewidth}}}]{
                \texttt{\tiny{[P2$\rightarrow$P1]}}
                \texttt{Post{-}episode Reflexion (not sent to GM):} \\
\texttt{New plan: In my last attempt, I failed to provide a description that would lead the other player to guess the target word in 3 tries. I started with a good description, but then I failed to provide a description that would help the other player to make a connection between the given clues and the target word. What I should have done is to provide a description that would give a hint about the word's meaning or its relation to the related words. However, I also failed to use the related words effectively. In my last attempt, I only used the related words as a starting point, but I didn't use them to create a connection between the clues. Next time, I will try to use the related words more effectively by creating a connection between the clues and the related words. For example, I could have said something like "It's a quality that is similar to being constant, but it's more related to movement or stability, like when you're walking on a firm surface." This would have given the other player a better chance to guess the target word. Additionally, I will try to use the given clues more effectively by creating a connection between them. For example, I could have said something like "Imagine a surface that doesn't wobble when you walk on it, and it's also a quality that helps you balance while moving, like when you're walking on a firm surface." This would have given the other player a clear hint about the target word's meaning.} \\
            }
        }
    }
    & & \\ \\

    \theutterance \stepcounter{utterance}  
    & & & \multicolumn{4}{p{0.6\linewidth}}{
        \cellcolor[rgb]{0.95,0.95,0.95}{
            \makecell[{{p{\linewidth}}}]{
                \texttt{\tiny{[P1$\rightarrow$P2]}}
                \texttt{Post{-}episode Reflexion (not sent to GM):} \\
\texttt{New plan: } \\
\\ 
\texttt{In my last attempt, I was able to guess two words that seemed related to the given clues, but I failed to guess the target word in three tries. Upon reflection, I realize that I was focusing too much on the literal meaning of the clues and not enough on the connections between them. } \\
\\ 
\texttt{Next time, I will try to find a word that ties together the three clues in a more abstract or metaphorical way. For example, I might look for a word that represents a quality that helps you walk on a tightrope (Balance) and is also a characteristic of a calm sea (Calmness), and is something that doesn't change or waver (Constant). } \\
\\ 
\texttt{I will also try to think more creatively and consider words that might not be the most obvious choices, but could still fit the given clues. To achieve this, I will take a step back and re{-}read the clues, looking for common themes or patterns that might not be immediately apparent. I will also try to think of words that are related to the clues in multiple ways, such as through synonyms, antonyms, or associations. } \\
\\ 
\texttt{Specifically, in this case, I will consider words that are related to the concept of stability or steadiness, and see if any of them fit the given clues. I will also think about words that are related to the idea of movement or balance, and see if any of them tie together the three clues in a meaningful way.} \\
\\ 
\texttt{To further improve my chances, I will also try to make connections between the clues and my own experiences or knowledge. For example, I might think about times when I had to balance on a tightrope or walk on a calm sea, and see if any words come to mind that relate to those experiences. By making these connections, I hope to come up with a word that is not only related to the clues, but also meaningful and relevant to my own life.} \\
            }
        }
    }
     \\ \\

\end{supertabular}
}

\end{document}
